% =====================================================================
%  סיכומי הרצאות — קורס 9921
%  קומפילציה: xelatex lecture_notes.tex  (פעמיים, לצורך תוכן עניינים והפניות)
% =====================================================================
\documentclass[11pt,a4paper]{report}

% ---------------------------------------------------------------------
%  מתגים: הצגה / הסתרה של הצעות ושל הערות "דורש התייחסות"
%  (להדפסה סופית — לשנות ל-false)
% ---------------------------------------------------------------------
\newif\ifshowsuggestions \showsuggestionstrue
\newif\ifshowtodos       \showtodostrue

% ---------------------------------------------------------------------
%  חבילות (tikz, hyperref ו-tcolorbox לפני polyglossia/bidi)
% ---------------------------------------------------------------------
\usepackage{amsmath,amssymb,amsthm,mathtools}
\usepackage[a4paper,margin=2.5cm,headheight=14pt]{geometry}
\usepackage{xcolor}
\usepackage{etoolbox}
\usepackage{tikz}
\usetikzlibrary{calc}
\usepackage[most]{tcolorbox}
\usepackage{enumitem}
\usepackage{array,booktabs}
\usepackage{titlesec}
\usepackage{fancyhdr}
\usepackage{comment}
\usepackage[hidelinks]{hyperref}
\usepackage{polyglossia}
\setmainlanguage{hebrew}
\setotherlanguage{english}

\setmainfont{Latin Modern Roman}
\newfontfamily\hebrewfont[Script=Hebrew]{David CLM}
\newfontfamily\hebrewfontsf[Script=Hebrew]{Simple CLM}
\newfontfamily\hebrewfonttt[Script=Hebrew]{Miriam Mono CLM}

% ---------------------------------------------------------------------
%  צבעים
% ---------------------------------------------------------------------
\definecolor{defc}{RGB}{31,78,121}      % הגדרות
\definecolor{thmc}{RGB}{20,110,100}     % טענות ומשפטים
\definecolor{exc}{RGB}{110,110,110}     % דוגמאות
\definecolor{keyc}{RGB}{150,40,40}      % נוסחאות חשובות
\definecolor{sugc}{RGB}{214,51,132}     % הצעות (ורוד)
\definecolor{todoc}{RGB}{200,110,0}     % דורש התייחסות (כתום)

% ---------------------------------------------------------------------
%  כותרות ועימוד
% ---------------------------------------------------------------------
\titleformat{\chapter}[display]
  {\sffamily\bfseries\color{defc}}
  {\Large\ifnum\value{chapter}=0 פרק מקדים\else פרק \thechapter\fi}
  {6pt}
  {\Huge}
  [\vspace{2pt}{\color{defc}\titlerule[1pt]}]
\titlespacing*{\chapter}{0pt}{-10pt}{24pt}
\titleformat{\section}{\sffamily\Large\bfseries\color{defc}}{\thesection}{0.8em}{}
\titleformat{\subsection}{\sffamily\large\bfseries\color{defc!85!black}}{\thesubsection}{0.8em}{}

\pagestyle{fancy}
\fancyhf{}
\fancyhead[R]{\small\sffamily\leftmark}
\fancyhead[L]{\small\sffamily קורס 9921}
\fancyfoot[C]{\small\thepage}
\renewcommand{\headrulewidth}{0.4pt}
\renewcommand{\chaptermark}[1]{\markboth{\ifnum\value{chapter}=0 פרק מקדים\else פרק \thechapter\fi: #1}{}}
\fancypagestyle{plain}{\fancyhf{}\fancyfoot[C]{\small\thepage}\renewcommand{\headrulewidth}{0pt}}

\setlist{itemsep=2pt,topsep=4pt}
\setlist[itemize]{label=\small$\bullet$}
\setlength{\parindent}{0pt}
\setlength{\parskip}{5pt}

% ---------------------------------------------------------------------
%  סביבות מתמטיות (מספור משותף בתוך כל פרק)
% ---------------------------------------------------------------------
\tcbset{
  base/.style={enhanced, breakable, sharp corners=south, arc=2pt,
    fonttitle=\sffamily\bfseries, before skip=10pt, after skip=10pt,
    left=6pt, right=6pt, top=4pt, bottom=4pt,
    separator sign={:}, description delimiters none},
}

\newtcbtheorem[number within=chapter]{definition}{הגדרה}{base,
  colback=defc!4, colframe=defc, coltitle=white}{def}

\newtcbtheorem[use counter from=definition]{claim}{טענה}{base,
  colback=thmc!4, colframe=thmc, coltitle=white}{clm}

\newtcbtheorem[use counter from=definition]{principle}{עקרון}{base,
  colback=thmc!4, colframe=thmc, coltitle=white}{prn}

\newtcbtheorem[use counter from=definition]{theorem}{משפט}{base,
  colback=thmc!4, colframe=thmc, coltitle=white}{thm}

\newtcbtheorem[use counter from=definition]{corollary}{מסקנה}{base,
  colback=thmc!4, colframe=thmc, coltitle=white}{cor}

\newtcbtheorem[use counter from=definition]{example}{דוגמה}{base,
  colback=white, colframe=exc!60, coltitle=exc!20!black, colbacktitle=exc!12,
  boxrule=0.5pt}{ex}

% תיבת נוסחה / סיכום חשוב
\newtcolorbox{keybox}[1][]{enhanced, breakable, colback=keyc!5, colframe=keyc!80,
  boxrule=0.8pt, arc=2pt, left=8pt, right=8pt, before skip=8pt, after skip=8pt,
  fonttitle=\sffamily\bfseries, coltitle=white, #1}

% הערה קצרה (ללא מסגרת)
\newenvironment{remark}{\par\smallskip\noindent\textsf{\textbf{הערה.}}\ }{\par\smallskip}
\newenvironment{notation}{\par\smallskip\noindent\textsf{\textbf{סימון.}}\ }{\par\smallskip}

% הוכחה
\renewcommand{\proofname}{הוכחה}

% ---------------------------------------------------------------------
%  הצעות לתוספות (ורוד)
%  שימוש: \begin{suggestion}{דוגמה נוספת} ... \end{suggestion}
% ---------------------------------------------------------------------
\newtcolorbox{suggestion}[1]{enhanced, breakable, colback=sugc!6, colframe=sugc,
  boxrule=0.6pt, arc=3pt, left=8pt, right=8pt, top=4pt, bottom=4pt,
  before skip=10pt, after skip=10pt,
  attach boxed title to top right={xshift=-8pt, yshift=-7pt},
  boxed title style={colback=sugc, colframe=sugc, arc=2pt},
  fonttitle=\sffamily\bfseries\small, coltitle=white,
  title={הצעה: #1}}

% ---------------------------------------------------------------------
%  "דורש התייחסות" (כתום) — נאסף אוטומטית לרשימה בתחילת המסמך
%  \needcheck{טקסט}           — הערה קצרה בתוך השורה
%  \begin{attention}{תקציר} ... \end{attention} — תיבה
% ---------------------------------------------------------------------
% דגל קטן (הגופנים העבריים אינם כוללים את התו)
\newcommand{\flag}{\tikz[baseline=-0.1ex,x=0.9ex,y=0.9ex]{\draw[line width=0.5pt] (0,-0.3)--(0,1.6); \fill (0,1.6)--(1.2,1.2)--(0,0.8)--cycle;}}
\newcounter{todo}
\makeatletter
\newcommand*\l@todo{\@dottedtocline{1}{0em}{2.2em}}
\newcommand{\listoftodos}{\@starttoc{tdo}}
\makeatother
\newcommand{\addtodo}[1]{\refstepcounter{todo}%
  \addcontentsline{tdo}{todo}{\protect\numberline{\thetodo}#1}}

\newcommand{\needcheck}[1]{\addtodo{#1}%
  {\color{todoc}\sffamily\small\,[\textbf{\flag{}\,\thetodo\ לבדיקה:} #1]\,}}

\tcbset{attstyle/.style={enhanced, breakable, colback=todoc!7, colframe=todoc,
  boxrule=0.6pt, arc=2pt, left=8pt, right=8pt, before skip=10pt, after skip=10pt,
  borderline={0.6pt}{2pt}{todoc, dashed},
  fonttitle=\sffamily\bfseries\small, coltitle=white}}
\newtcolorbox{attention}[1]{attstyle, before upper={\addtodo{#1}},
  title={דורש התייחסות: #1}}
\newtcolorbox{attentiondemo}[1]{attstyle, title={דורש התייחסות: #1}}

\ifshowsuggestions\else\excludecomment{suggestion}\fi
\ifshowtodos\else
  \excludecomment{attention}
  \excludecomment{attentiondemo}
  \renewcommand{\needcheck}[1]{}
\fi

% ---------------------------------------------------------------------
%  פקודות מתמטיות
% ---------------------------------------------------------------------
\newcommand{\N}{\mathbb{N}}
\newcommand{\Z}{\mathbb{Z}}
\newcommand{\Q}{\mathbb{Q}}
\newcommand{\R}{\mathbb{R}}
\newcommand{\set}[1]{\left\{#1\right\}}
\newcommand{\setof}[2]{\left\{\,#1 \;\middle|\; #2\,\right\}}
\newcommand{\symdiff}{\mathbin{\triangle}}
\newcommand{\comp}[1]{\overline{#1}}
\newcommand{\Pow}{\mathcal{P}}
\newcommand{\heb}[1]{\text{\RL{#1}}}
\renewcommand{\emptyset}{\varnothing}

% קו מפריד בכותרות (לגופן Simple CLM אין קו מפריד ארוך)
\newcommand{\dash}{{\rmfamily —}}

% סימון תחילת הרצאה
\newcommand{\lecture}[1]{\par\medskip
  \begin{center}\color{exc}\sffamily\small
  \rule[0.5ex]{0.25\linewidth}{0.4pt}\quad הרצאה #1\quad\rule[0.5ex]{0.25\linewidth}{0.4pt}
  \end{center}\par}

% =====================================================================
\begin{document}
% =====================================================================

\begin{titlepage}
  \centering
  \vspace*{3cm}
  {\sffamily\color{defc}\rule{\linewidth}{1pt}\par\vspace{10pt}
   {\Huge\bfseries סיכומי הרצאות\par}\vspace{8pt}
   {\LARGE קורס 9921\par}\vspace{10pt}
   \rule{\linewidth}{1pt}\par}
  \vspace{1.5cm}
  {\Large מרצה: שובל\par}
  \vspace{0.5cm}
  \ifshowtodos
  \begin{minipage}{0.8\linewidth}\centering
  \needcheck{להשלים את שם הקורס, שם המרצה המלא והסמסטר}
  \end{minipage}
  \fi
  \vfill
  {\large נכתב על בסיס רשימות שנכתבו בהרצאות\par}
  \vspace{1cm}
\end{titlepage}

% ---------------------------------------------------------------------
\chapter*{על המסמך}
\addcontentsline{toc}{chapter}{על המסמך}

המסמך נבנה מתוך רשימות בכתב יד שנכתבו במהלך ההרצאות. התוכן נשמר כמעט כפי שהוא, עם סידור, ניסוח מחדש ותיקון של שגיאות כתיבה ברורות. לאורך המסמך מופיעים כמה סוגי תיבות:

\begin{definition*}{דוגמה לתיבה}
הגדרות (ובצבע טורקיז: טענות, עקרונות ומשפטים) — התוכן המרכזי של הקורס.
\end{definition*}

\begin{suggestion}{סוג ההצעה}
תיבה ורודה היא \textbf{הצעה} לתוספת שאינה מופיעה ברשימות המקוריות: דוגמה נוספת, הסבר נוסף, איור וכדומה. אפשר לאמץ, לערוך או למחוק.
\end{suggestion}

\begin{attentiondemo}{דוגמה לתיבה כתומה}
תיבה כתומה, או סימון כתום בתוך השורה \textcolor{todoc}{\sffamily\small[\flag{} לבדיקה: ...]}, מסמנת מקום שבו היה ספק בפענוח הרשימות, או מקום שבו תיקנתי משהו ויש לוודא שהתיקון נכון. כל המקומות האלה מרוכזים ברשימה שבהמשך העמוד.
\end{attentiondemo}

את שני סוגי התוספות אפשר להסתיר בבת אחת באמצעות המתגים \texttt{showsuggestions} ו-\texttt{showtodos} שבראש קובץ המקור.

\ifshowtodos
\section*{רשימת המקומות הדורשים התייחסות}
\listoftodos
\fi

\tableofcontents

% =====================================================================
\setcounter{chapter}{-1}
\chapter{תורת הקבוצות}
% =====================================================================

\lecture{1}

\section{מהי קבוצה?}

\begin{definition}{קבוצה~\dash{} לא פורמלית}{set}
\textbf{קבוצה} היא אוסף של \textbf{איברים}, כאשר:
\begin{itemize}
  \item אין חשיבות לסדר האיברים;
  \item אין משמעות לחזרה על איברים.
\end{itemize}
\end{definition}

\begin{example}{}{}
\[
  \set{1,2,3} \;=\; \set{3,2,1} \;=\; \set{1,1,2,3,1}.
\]
\end{example}

\begin{notation}
אם $x$ שייך לקבוצה $A$, נסמן $x\in A$, כלומר $x$ \emph{איבר} ב-$A$.
בדומה, $x\notin A$ פירושו ש-$x$ אינו איבר ב-$A$.
\end{notation}

\section{קבוצות של מספרים}

נסמן כמה קבוצות מספרים מוכרות:
\begin{align*}
  \N &= \set{1,2,3,\dots} && \heb{המספרים הטבעיים}\\
  \Z &= \set{\dots,-2,-1,0,1,2,\dots} && \heb{המספרים השלמים}\\
  \Q &= \setof{\tfrac{m}{n}}{m,n\in\Z,\ n\neq 0} && \heb{המספרים הרציונליים}\\
  \R & && \heb{המספרים הממשיים}
\end{align*}
\needcheck{בהגדרת $\Q$ ברשימות נכתב $m\neq 0$; תיקנתי ל-$n\neq 0$ (המכנה הוא שאסור שיתאפס)}

למשל, $\frac13\in\Q$ ו-$-\frac{10}{4}\in\Q$.

\section{תיאור קבוצה על ידי תנאי}

באופן כללי, קבוצה נכתבת בצורה
\[
  \Big\{\ \underbrace{\square}_{\heb{מי שבקבוצה}}\ \ \underbrace{\Big|}_{\heb{"כך ש..."}}\ \ \underbrace{\triangle}_{\heb{התנאי}}\ \Big\}
\]

\begin{example}{}{}
\[
  A=\setof{2n}{n\in\N}.
\]
במילים: $A$ מכילה את כל הכפולות ב-$2$ של המספרים הטבעיים, כלומר $A=\set{2,4,6,\dots}$.
\end{example}

\begin{example}{}{}
\[
  B=\setof{x}{1<x<7,\ x\in\N}=\set{2,3,4,5,6}.
\]
הפסיק בתוך התנאי פירושו \emph{"וגם"}: $x$ מקיים $1<x<7$ \emph{וגם} $x\in\N$.
\end{example}

\section{יחס ההכלה}

\begin{definition}{הכלה}{subset}
יהיו $A,B$ קבוצות. נאמר כי $A$ \textbf{מוכלת} ב-$B$, או ש-$A$ היא \textbf{תת-קבוצה} של $B$, ונסמן $A\subseteq B$, אם כל איבר ב-$A$ נמצא גם ב-$B$.
\end{definition}

\begin{example}{}{}
$\N\subseteq\Z$, כי $\set{1,2,3,\dots}\subseteq\set{\dots,-1,0,1,2,\dots}$.

\smallskip
לעומת זאת, $\Z\nsubseteq\N$, כי קיים איבר ב-$\Z$ שאינו ב-$\N$. למשל $-1$, או $0$.
\end{example}

\begin{definition}{הקבוצה הריקה}{empty}
\textbf{הקבוצה הריקה} היא הקבוצה שאין בה אף איבר. נסמן אותה $\emptyset$, ומתקיים: לכל $x$, $x\notin\emptyset$.
באופן שקול: לא קיים $x$ שעבורו $x\in\emptyset$.
\end{definition}

\section{פעולות על קבוצות}

\begin{definition}{פעולות על קבוצות}{ops}
יהיו $A,B$ קבוצות.
\begin{enumerate}
  \item \textbf{חיתוך:}\quad $A\cap B=\setof{x}{x\in A,\ x\in B}$.
  \item \textbf{איחוד:}\quad $A\cup B=\setof{x}{x\in A \ \heb{או}\ x\in B}$.
  \item \textbf{חיסור / הפרש} ("$A$ פחות $B$", כלומר "מי שב-$A$ ולא ב-$B$"):\quad $A\setminus B=\setof{x}{x\in A,\ x\notin B}$.
  \item \textbf{הפרש סימטרי} (מי שנמצא \emph{בדיוק} באחת מהקבוצות):\quad $A\symdiff B=(A\setminus B)\cup(B\setminus A)$.
\end{enumerate}
\end{definition}

\begin{example}{}{}
יהיו $A=\set{1,2,3}$ ו-$B=\set{2,3,4}$. אז
\[
\begin{aligned}
  A\setminus B &= \set{1}, & B\setminus A &= \set{4},\\
  A\cap B &= \set{2,3}, & A\cup B &= \set{1,2,3,4},\\
  A\symdiff B &= \set{1,4}.
\end{aligned}
\]
\end{example}

\begin{suggestion}{איור~\dash{} דיאגרמות ון}
דיאגרמות ון עוזרות לזכור את ההגדרות (האזור הצבוע הוא תוצאת הפעולה):

\medskip
\begin{center}
\newcommand{\venn}[2]{%
\begin{tikzpicture}[scale=0.55]
  \def\A{(0,0) circle (1.2)}\def\B{(1.4,0) circle (1.2)}
  #1
  \draw[thick] \A; \draw[thick] \B;
  \node at (-0.55,0) {$A$}; \node at (1.95,0) {$B$};
  \node[below] at (0.7,-1.35) {#2};
\end{tikzpicture}}
\venn{\begin{scope}[even odd rule]\fill[sugc!35]\A \B;\end{scope}}{$A\symdiff B$}\hfill
\venn{\begin{scope}\begin{scope}[even odd rule]\clip\B (-2,-2) rectangle (4,2);\fill[sugc!35]\A;\end{scope}\end{scope}}{$A\setminus B$}\hfill
\venn{\fill[sugc!35]\A;\fill[sugc!35]\B;}{$A\cup B$}\hfill
\venn{\begin{scope}\clip\A;\fill[sugc!35]\B;\end{scope}}{$A\cap B$}
\end{center}
\end{suggestion}

\begin{claim}{}{basic}
לכל שתי קבוצות $A,B$ מתקיים:
\[
\renewcommand{\arraystretch}{1.3}
\begin{array}{lll}
  (1)\ \ A\cap B\subseteq A & (4)\ \ A\subseteq A\cup B & (7)\ \ A\setminus B\subseteq A\\
  (2)\ \ A\cap B\subseteq B & (5)\ \ B\subseteq A\cup B & (8)\ \ A\setminus B\subseteq A\symdiff B\\
  (3)\ \ A\cap B= B\cap A & (6)\ \ A\cup B= B\cup A & (9)\ \ A\symdiff B= B\symdiff A
\end{array}
\]
\end{claim}

\begin{suggestion}{דוגמה להוכחה}
ברשימות הטענה מופיעה ללא הוכחה. כדאי להראות לפחות סעיף אחד, כדי להדגים את השיטה הסטנדרטית להוכחת הכלה: \emph{לוקחים איבר כלשהו בצד שמאל ומראים שהוא שייך לצד ימין.}

\textbf{הוכחת (1).} יהי $x\in A\cap B$. לפי הגדרת החיתוך $x\in A$ וגם $x\in B$; בפרט $x\in A$. מכיוון ש-$x$ היה איבר כלשהו של $A\cap B$, קיבלנו $A\cap B\subseteq A$.

\textbf{הוכחת (8).} יהי $x\in A\setminus B$. מאחר ש-$A\symdiff B=(A\setminus B)\cup(B\setminus A)$, ו-$x$ שייך לקבוצה הראשונה באיחוד, נקבל $x\in A\symdiff B$.

\smallskip
כדי להוכיח \emph{שוויון} $X=Y$ מראים שתי הכלות: $X\subseteq Y$ וגם $Y\subseteq X$.
\end{suggestion}

\subsection{המשלים}

\begin{definition}{משלים}{complement}
תהי $U$ קבוצה "אוניברסלית" (כל האיברים שבהם עוסקים), ותהי $A\subseteq U$. \textbf{המשלים} של $A$, שיסומן $\comp{A}$ (או $A^c$), הוא
\[
  \comp{A}=A^c=\setof{x\in U}{x\notin A}=U\setminus A.
\]
\end{definition}
\needcheck{ברשימות ההגדרה נכתבה בלי קבוצה אוניברסלית, $\comp{A}=\setof{x}{x\notin A}$, ורק בדוגמה הופיע $U=\N$. הוספתי את $U$ להגדרה. יש לוודא שזה תואם את ההגדרה של המרצה}

\begin{example}{$U=\N$}{}
\begin{itemize}
  \item $A=\set{1,2,3}$, ולכן $\comp{A}=\set{4,5,6,\dots}=\N\setminus A$.
  \item $A=\N\setminus\set{2,4}$, ולכן $\comp{A}=\set{2,4}$.
  \item $A=\N$, ולכן $\comp{A}=\emptyset$.
\end{itemize}
\end{example}

\begin{suggestion}{תוספת תוכן~\dash{} כללי דה־מורגן}
טבעי להוסיף כאן את כללי דה־מורגן, שמקשרים בין משלים, איחוד וחיתוך:
\[
  \comp{A\cup B}=\comp{A}\cap\comp{B},\qquad \comp{A\cap B}=\comp{A}\cup\comp{B}.
\]
\emph{בדיקה על דוגמה} ($U=\N$, $A=\set{1,2,3}$, $B=\set{2,3,4}$):
$\comp{A\cup B}=\set{5,6,\dots}$, ומצד שני $\comp{A}\cap\comp{B}=\set{4,5,\dots}\cap\set{1,5,6,\dots}=\set{5,6,\dots}$.
\end{suggestion}

\section{קבוצת החזקה}

\begin{definition}{קבוצת החזקה}{power}
לכל קבוצה $A$ נסמן ב-$\Pow(A)$, או ב-$2^A$, את הקבוצה שאיבריה הם \emph{כל} תתי-הקבוצות של $A$:
\[
  \Pow(A)=\setof{B}{B\subseteq A}.
\]
\end{definition}

\begin{example}{}{}
עבור $A=\set{1,2}$, תתי-הקבוצות של $A$ הן $\emptyset,\ \set{1},\ \set{2},\ \set{1,2}$, ולכן
\[
  \Pow(A)=\big\{\emptyset,\ \set{1},\ \set{2},\ \set{1,2}\big\}.
\]
כאן יש ב-$\Pow(A)$ ארבעה איברים.
\end{example}

\begin{example}{}{powabc}
עבור $A=\set{a,b,c}$, תתי-הקבוצות הן:
\[
  \emptyset,\quad \set{a},\ \set{b},\ \set{c},\quad \set{a,b},\ \set{a,c},\ \set{b,c},\quad \set{a,b,c}.
\]
סך הכול $8=2^3$ תתי-קבוצות, ול-$A$ יש $3$ איברים.
\needcheck{ברשימות $\set{a,c}$ הופיעה פעמיים; תיקנתי את המופע השני ל-$\set{b,c}$}
\end{example}

\begin{remark}
לכל קבוצה $A$ תמיד מתקיים $\emptyset\subseteq A$ וגם $A\subseteq A$ (כל קבוצה מוכלת בעצמה). לכן $\emptyset$ ו-$A$ הם תמיד איברים של $\Pow(A)$.
בהמשך (מסקנה~\ref{cor:powsize}) נראה שאם ל-$A$ יש $n$ איברים, אז ל-$\Pow(A)$ יש $2^n$ איברים.
\end{remark}

\section{זוג סדור ומכפלה קרטזית}

\begin{definition}{מכפלה קרטזית}{cartesian}
יהיו $A,B$ קבוצות. \textbf{המכפלה הקרטזית} של $A$ ו-$B$ היא \emph{הקבוצה}
\[
  A\times B=\setof{(a,b)}{a\in A,\ b\in B},
\]
כאשר $(a,b)$ הוא \textbf{זוג סדור}.
\end{definition}

\begin{example}{}{}
$A=\set{1,2}$, $B=\set{3,4}$. אז
\[
  A\times B=\set{(1,3),(1,4),(2,3),(2,4)}.
\]
מצד שני,
\[
  B\times A=\set{(3,1),(3,2),(4,1),(4,2)}.
\]
\textbf{שימו לב:} הזוגות הסדורים $(1,3)$ ו-$(3,1)$ \emph{שונים}! בזוג סדור הסדר חשוב, בניגוד לקבוצה.
\end{example}

\begin{example}{}{}
מצאו את כל האיברים בקבוצה $\Pow(\set{1})\times\Pow(\emptyset)$.
\[
  \Pow(\set{1})\times\Pow(\emptyset)
  =\underbrace{\set{\emptyset,\set{1}}}_{A}\times\underbrace{\set{\emptyset}}_{B}
  =\big\{(\emptyset,\emptyset),\ (\set{1},\emptyset)\big\}.
\]
\end{example}

\begin{suggestion}{הסבר נוסף}
נקודה עדינה בדוגמה האחרונה: $\Pow(\emptyset)=\set{\emptyset}$ היא קבוצה עם \emph{איבר אחד} (הקבוצה הריקה), ולכן היא אינה ריקה. אילו היינו כופלים ב-$\emptyset$ עצמה, היינו מקבלים $A\times\emptyset=\emptyset$.
כדאי לשים לב גם לגדלים: $|A\times B|=|A|\cdot|B|$ (כאן $2\cdot1=2$). זו הצצה ל\emph{עקרון הכפל} שבפרק הבא.
\end{suggestion}

% =====================================================================
\chapter{קומבינטוריקה}
% =====================================================================

\section{מבוא: ספירה}

קומבינטוריקה עוסקת במנייה ובספירה של דברים. נרצה לענות על שאלות מהצורה \emph{"בכמה דרכים ניתן...?"}, ולשם כך נכיר עקרונות כלליים ונוסחאות. נפתח בשלוש שאלות לדוגמה, ונשים לב למאפיינים של כל אחת:

\begin{enumerate}[label=\textbf{(\arabic*)}]
  \item \textbf{בכמה דרכים ניתן לבחור 3 סטודנטים מתוך 100?}
  \begin{itemize}
    \item הבחירה היא \emph{ללא חזרות}: לא בוחרים את אותו סטודנט פעמיים.
    \item \emph{אין חשיבות לסדר} הבחירה.
  \end{itemize}

  \item \textbf{לאורך הסמסטר יש 3 מטלות, ולכל מטלה יש בודק (אחד מתוך 100 הסטודנטים). כמה הרכבים של בודקים יש?}
  \begin{itemize}
    \item \emph{יש אפשרות לחזרות}: למשל, יוסי יבדוק את כל המטלות.
    \item \emph{יש חשיבות לסדר}: יש הבדל בין מי שבודק את מטלה 1 לבין מי שבודק את מטלה 2.
  \end{itemize}

  \item \textbf{מנעול עם קוד בן 3 ספרות (כל ספרה בין 0 ל-9). כמה קודים אפשריים יש?}
  \begin{itemize}
    \item \emph{יש חשיבות לסדר.}
    \item \emph{יש חזרות.}
  \end{itemize}
\end{enumerate}

\begin{suggestion}{לסגור את המעגל}
כדאי לחזור לשלוש השאלות אחרי סעיף~\ref{sec:choose} ולתת להן תשובות:
\[
  (1)\ \binom{100}{3}=\frac{100\cdot 99\cdot 98}{3!}=161{,}700,\qquad
  (2)\ 100^3=1{,}000{,}000,\qquad
  (3)\ 10^3=1{,}000.
\]
\end{suggestion}

\lecture{2}

\section{עקרון הכפל}

\begin{principle}{עקרון הכפל}{mult}
אם נרצה לבצע משימה שמורכבת מ-$k$ שלבים, ובשלב ה-$i$ ניתן לבחור אחת מתוך $n_i$ אפשרויות, אז מספר הדרכים השונות לבצע את המשימה הוא
\[
  n_1\cdot n_2\cdot n_3\cdots n_k.
\]
\textbf{דגש:} מספר האפשרויות בכל שלב אינו תלוי במה שנבחר בשלב אחר.
\end{principle}
\needcheck{ברשימות נכתב "אחת מתוך $n_k$", ותיקנתי ל-$n_i$. כמו כן, יש לוודא את פענוח פריטי הארוחה בדוגמה הבאה (מיצים / סלטים / ביצים / לחמים)}

\begin{example}{ארוחת בוקר}{}
אני מרכיב ארוחת בוקר, כאשר ניתן לבחור:
\begin{itemize}
  \item אחד מתוך 3 מיצים \hfill ($n_1=3$)
  \item אחד מתוך 4 סלטים \hfill ($n_2=4$)
  \item אחד מתוך 3 סוגי ביצים \hfill ($n_3=3$)
  \item אחד מתוך 2 סוגי לחמים \hfill ($n_4=2$)
\end{itemize}
יש כאן $k=4$ שלבים, ולכן מספר ארוחות הבוקר השונות שניתן להרכיב הוא
$3\cdot4\cdot3\cdot2=72$.
\end{example}

\begin{example}{הושבה בשורה}{seating}
מושיבים 10 אנשים בשורה של 10 כיסאות. בכמה דרכים שונות אפשר להושיב אותם?
\begin{itemize}
  \item המשימה הראשונה היא להושיב מישהו בכיסא 1: \quad $n_1=10$.
  \item המשימה השנייה היא להושיב מישהו בכיסא 2: \quad $n_2=9$.
  \item[] \hfill$\vdots$\hfill\null
  \item המשימה העשירית היא להושיב מישהו בכיסא 10: \quad $n_{10}=1$.
\end{itemize}
לכן מספר הדרכים הוא
\[
  10\cdot9\cdot8\cdot7\cdots1=10!.
\]
\end{example}

\begin{suggestion}{הסבר נוסף~\dash{} עצרת, ומה בדיוק אומר ה"דגש"}
\textbf{עצרת.} לכל $n\in\N$ מגדירים $n!=1\cdot2\cdots n$, ובנוסף $0!=1$ (נשתמש בזה בהמשך). הסימון מופיע בדוגמה בלי הגדרה, אז כדאי להגדיר אותו במפורש.

\textbf{הדגש בעקרון הכפל.} בדוגמת ההושבה, \emph{מי} יכול לשבת בכיסא 2 תלוי במי שהושב בכיסא 1, אבל \emph{מספר} האפשרויות הוא תמיד 9. זה כל מה שעקרון הכפל דורש: שמספר האפשרויות בכל שלב יהיה קבוע, גם אם האפשרויות עצמן משתנות.
\end{suggestion}

\section{עקרון החיבור}

\begin{notation}
לקבוצה $A$ נסמן ב-$|A|$ את \textbf{גודל} הקבוצה, כלומר את מספר האיברים בקבוצה.
\end{notation}

\begin{definition}{קבוצות זרות}{disjoint}
הקבוצות $A,B$ נקראות \textbf{זרות} אם $A\cap B=\emptyset$, כלומר אין להן אף איבר משותף.
\end{definition}

\begin{principle}{עקרון החיבור~\dash{} שתי קבוצות}{add2}
אם $A,B$ קבוצות זרות, אז
\[
  |A\cup B|=|A|+|B|.
\]
\end{principle}

\begin{example}{}{}
יהיו $A=\set{1,2}$, $B=\set{3,4}$ ו-$C=\set{2,3,4}$.
\begin{itemize}
  \item $A,B$ זרות, ולכן $A\cup B=\set{1,2,3,4}$ ואכן $|A\cup B|=4=|A|+|B|=2+2$.
  \item $A,C$ \emph{אינן} זרות (שתיהן מכילות את $2$). כאן $A\cup C=\set{1,2,3,4}$, ולכן
  \[|A\cup C|=4\neq 5=|A|+|C|.\]
\end{itemize}
\end{example}

\begin{definition}{חלוקה}{partition}
תהי $A$ קבוצה. \textbf{חלוקה} של $A$ היא אוסף של קבוצות $A_1,\dots,A_n$ המוכלות ב-$A$, שמקיימות:
\begin{itemize}
  \item $A=A_1\cup A_2\cup\dots\cup A_n$, וגם
  \item $A_i\cap A_j=\emptyset$ לכל $i\neq j$ ב-$\set{1,\dots,n}$.
\end{itemize}
\end{definition}

\begin{example}{}{partex}
ניקח $A=\set{1,2,3,4}$. האם הקבוצות הבאות מהוות חלוקה של $A$?
\begin{enumerate}[label=(\alph*)]
  \item $A_1=\set{1,2}$, $A_2=\set{3}$. \quad \textbf{לא}, כי $A_1\cup A_2\neq A$ (חסר $4$).
  \item $A_1=\set{1}$, $A_2=\set{2,3}$, $A_3=\set{4}$. \quad \textbf{כן!} $A_1\cup A_2\cup A_3=A$ והקבוצות זרות בזוגות.
  \item $A_1=\set{1,2,3}$, $A_2=\set{3,4}$. \quad \textbf{לא}, כי $A_1\cap A_2\neq\emptyset$.
  \item $A_1=\emptyset$, $A_2=\set{1,2,3,4}$. \quad \textbf{כן!}
  \item $A_1=\set{1,2,3}$, $A_2=\set{4,5}$. \quad \textbf{לא!} $A_2$ אינה מוכלת ב-$A$.
\end{enumerate}
\needcheck{בסעיפים (א) ו-(ב) נכתב ברשימות $A_1\cap A_2\neq A$ ו-$A_1\cap A_2\cap A_3=A$; תיקנתי ל-$\cup$ (איחוד)}\par
\needcheck{בסעיף (ד) הקבוצה הריקה מתקבלת כחלק בחלוקה. בספרים רבים דורשים שכל $A_i$ יהיה לא ריק. לוודא מה המוסכמה בקורס}
\end{example}

\begin{remark}
\textbf{הכלה מול שייכות.} $\emptyset\subseteq A$ נכון \emph{לכל} קבוצה $A$. לעומת זאת $\emptyset\in A$ \emph{לא} נכון לכל קבוצה. למשל $\emptyset\notin\set{1,2,3,4}$, אבל $\emptyset\in\Pow(A)$ לכל $A$.
\end{remark}

\begin{principle}{עקרון החיבור~\dash{} כללי}{add}
אם $A_1,\dots,A_n$ היא חלוקה של $A$, אז
\[
  |A|=|A_1|+|A_2|+\dots+|A_n|.
\]
\end{principle}

בדרך כלל, השימוש בעקרון הוא כך: מחלקים את הבעיה ל\emph{מקרים} זרים (כלומר מוצאים חלוקה), סופרים כל מקרה בנפרד ומחברים.

\begin{example}{כתיבת מספר כסכום}{compositions}
בכמה דרכים ניתן לכתוב את $3$ כסכום של מספרים טבעיים (לא כולל $0$), כאשר יש משמעות לסדר?
\[
  3=1+2,\qquad 3=2+1,\qquad 3=1+1+1,\qquad 3=3.
\]
התשובה היא $4$.

\medskip
\textbf{אותה שאלה עבור $4$.} תהי $A$ קבוצת כל הדרכים לכתוב את $4$ כסכום; אנו מחפשים את $|A|$.
לכל $k=1,2,3,4$ נגדיר $A_k$ = קבוצת הדרכים לכתוב את $4$ כסכום \emph{שמתחיל} ב-$k$.
אז $A_1,\dots,A_4$ היא חלוקה של $A$:
\begin{center}
\renewcommand{\arraystretch}{1.2}
\begin{tabular}{c|c|c|c}
  $A_1$ & $A_2$ & $A_3$ & $A_4$\\ \hline
  $1+1+1+1$ & $2+1+1$ & $3+1$ & $4$\\
  $1+1+2$   & $2+2$   &       &    \\
  $1+2+1$   &         &       &    \\
  $1+3$     &         &       &
\end{tabular}
\end{center}
ולכן
\[
  |A|=|A_1|+|A_2|+|A_3|+|A_4|=4+2+1+1=8.
\]
\textbf{באופן כללי:} את $n$ ניתן לכתוב כסכום (עם משמעות לסדר) ב-$2^{n-1}$ דרכים.
\end{example}

\begin{suggestion}{הסבר נוסף~\dash{} למה $2^{n-1}$?}
ברשימות הנוסחה מופיעה בלי הסבר. הנה הוכחה קצרה בעזרת עקרון הכפל: נכתוב את $n$ כשורה של $n$ אחדות,
\[
  1\ \_\ 1\ \_\ 1\ \_\ \cdots\ \_\ 1 ,
\]
שביניהן $n-1$ רווחים. בכל רווח מחליטים אם לשים "$+$" (מפריד בין מחוברים) או להשאיר ריק (כלומר לאחד את האחדות). למשל עבור $n=4$, הבחירה $1+1\,1+1$ מתאימה לסכום $1+2+1$.
לכל רווח יש $2$ אפשרויות, ולכן יש $2^{n-1}$ סכומים שונים. עבור $n=4$ מתקבל $2^3=8$, בהתאם לחישוב.
\end{suggestion}

\section{בחירה של \texorpdfstring{$k$}{k} מתוך \texorpdfstring{$n$}{n} עצמים}\label{sec:choose}

בוחרים $k$ עצמים מתוך $n$. יש ארבעה מקרים, לפי שתי שאלות: האם יש \emph{משמעות לסדר}, והאם מותרות \emph{חזרות}.

\begin{center}
\renewcommand{\arraystretch}{1.6}
\begin{tabular}{r|c|c}
  & \textbf{עם חזרות} & \textbf{בלי חזרות}\\ \hline
  \textbf{יש משמעות לסדר} & $n^k$ \ \small(מקרה 1) & $\dfrac{n!}{(n-k)!}$ \ \small(מקרה 2)\\ \hline
  \textbf{אין משמעות לסדר} & {\color{todoc}?} \ \small(מקרה 4) & $\dbinom{n}{k}$ \ \small(מקרה 3)
\end{tabular}
\end{center}

\subsection*{מקרה 1: עם משמעות לסדר ועם חזרות}

\begin{example}{}{}
כל שבוע ממנים סטודנט למחוק את הלוח (סטודנט יכול להתמנות כמה פעמים). כמה שיבוצים אפשריים יש לסמסטר של 13 שבועות, כשיש 100 סטודנטים?
\[
  \underbrace{100}_{\heb{שבוע 1}}\cdot\underbrace{100}_{\heb{שבוע 2}}\cdots\underbrace{100}_{\heb{שבוע 13}}=100^{13}.
\]
יש $100^{13}$ דרכים לשיבוץ.
\end{example}

\begin{keybox}[title={באופן כללי}]
בחירה של $k$ מתוך $n$, \textbf{עם} משמעות לסדר ו\textbf{עם} חזרות: \qquad $n^k$.
\end{keybox}

\subsection*{מקרה 2: עם משמעות לסדר ובלי חזרות}

\begin{example}{}{}
כל שבוע בוחרים סטודנט למחוק את הלוח, אבל כל אחד ייבחר \emph{לכל היותר פעם אחת}. לפי עקרון הכפל:
\[
  \underbrace{100}_{\heb{שבוע 1}}\cdot\underbrace{99}_{\heb{שבוע 2}}\cdot 98\cdots\underbrace{88}_{\heb{שבוע 13}}.
\]
\end{example}

\begin{keybox}[title={באופן כללי}]
בחירה של $k$ מתוך $n$, \textbf{עם} משמעות לסדר ו\textbf{בלי} חזרות:
\[
  n(n-1)\cdots\underbrace{\big(n-(k-1)\big)}_{n-k+1}=\frac{n!}{(n-k)!}.
\]
\end{keybox}

\subsection*{מקרה 3: בלי משמעות לסדר ובלי חזרות}

\begin{example}{}{}
בכמה דרכים ניתן לבחור 2 סטודנטים (מתוך 100) שימחקו את הלוח?

נתחיל בדוגמה קטנה: בחירה של 2 מתוך 4, מתוך הקבוצה $\set{A,B,C,D}$:
\begin{center}
\renewcommand{\arraystretch}{1.25}
\begin{tabular}{c|c}
  \textbf{אין חזרות, אין סדר} & \textbf{אין חזרות, יש סדר}\\ \hline
  $\set{A,B}$ & $AB,\ BA$\\
  $\set{A,C}$ & $AC,\ CA$\\
  $\set{A,D}$ & $AD,\ DA$\\
  $\set{B,C}$ & $BC,\ CB$\\
  $\set{B,D}$ & $BD,\ DB$\\
  $\set{C,D}$ & $CD,\ DC$\\ \hline
  $6$ אפשרויות & $\dfrac{4!}{(4-2)!}=\dfrac{24}{2}=12$
\end{tabular}
\end{center}
\needcheck{ברשימות הופיע בטבלה $BD,\ BD$; תיקנתי ל-$BD,\ DB$}

ובחזרה לשאלה המקורית:
\[
  \binom{100}{2}=\frac{100!}{(100-2)!\cdot2!}=\frac{100!}{98!\cdot2}=\frac{100\cdot99}{2}=4950.
\]
\end{example}

\begin{keybox}[title={נוסחה~\dash{} המקדם הבינומי}]
בחירה של $k$ מתוך $n$, \textbf{בלי} משמעות לסדר ו\textbf{בלי} חזרות:
\[
  \binom{n}{k}=\frac{n!}{(n-k)!\,k!}.
\]
המספר $\binom{n}{k}$ נקרא \textbf{מקדם בינומי}, ונקרא "$n$ מעל $k$" (\LR{"$n$ choose $k$"}).
\end{keybox}

\begin{suggestion}{הסבר נוסף~\dash{} למה מחלקים ב-$k!$?}
הטבלה מראה את הרעיון: כל בחירה \emph{לא סדורה} $\set{A,B}$ מתאימה בדיוק ל-$2=2!$ בחירות \emph{סדורות} ($AB$ ו-$BA$). באופן כללי, כל קבוצה של $k$ עצמים אפשר לסדר ב-$k!$ דרכים (בדיוק כמו בדוגמת ההושבה~\ref{ex:seating}). לכן
\[
  \underbrace{\frac{n!}{(n-k)!}}_{\heb{עם סדר}}=\underbrace{\binom{n}{k}}_{\heb{בלי סדר}}\cdot\, k!
  \qquad\Longrightarrow\qquad
  \binom{n}{k}=\frac{n!}{(n-k)!\,k!}.
\]
\end{suggestion}

\subsection*{מקרה 4: בלי משמעות לסדר ועם חזרות}

\begin{attention}{המקרה הרביעי חסר ברשימות}
ברשימות מופיעים רק שלושת המקרים הראשונים. ייתכן שהמקרה הרביעי נלמד בהרצאה מאוחרת יותר, או שלא נלמד כלל. אם הוא נלמד, הנוסחה המקובלת היא $\binom{n+k-1}{k}$ (שיטת "כדורים ומחיצות").
\end{attention}

\section{תכונות פשוטות של \texorpdfstring{$\binom{n}{k}$}{(n choose k)}}

\begin{enumerate}[label=\textbf{\arabic*.}]
  \item \textbf{בכמה דרכים ניתן לבחור $n$ עצמים מתוך $n$ העצמים?} בדרך אחת בלבד, כלומר $\binom{n}{n}=1$. מצד שני, גם מהנוסחה:
  $\binom{n}{n}=\frac{n!}{n!\cdot0!}=1$.

  \item $\binom{n}{0}=1$. אפשר להגיע לזה מהנוסחה: $\binom{n}{0}=\frac{n!}{0!\cdot n!}=1$.

  \item \textbf{בכמה דרכים אפשר לבחור סטודנט אחד מתוך $n$?} $\binom{n}{1}=n$, ואכן
  $\binom{n}{1}=\frac{n!}{1!\cdot(n-1)!}=n$.\newline
  \needcheck{ברשימות נכתב בסוף החישוב $=n!$; תיקנתי ל-$=n$}

  \item \textbf{99 סטודנטים מתוך 100} (ובאופן כללי, $n-1$ מתוך $n$):
  $\binom{n}{n-1}=\frac{n!}{\big(n-(n-1)\big)!\cdot(n-1)!}=n$.

  \item \textbf{סימטריה:} $\displaystyle\binom{n}{k}=\binom{n}{n-k}$.

  \item \textbf{בכמה דרכים ניתן לבחור 3 סטודנטים מתוך 5?}
  $\binom{5}{3}=\frac{5!}{3!\cdot2!}=\frac{4\cdot5}{2}=10$.
\end{enumerate}

\begin{remark}
$\binom{n}{k}$ מוגדר (ובעל משמעות) רק כאשר $n,k$ שלמים ו-$0\le k\le n$.
\end{remark}

\section{הבינום של ניוטון}

נוסחאות מוכרות מהתיכון:
\[
  (a+b)^2=a^2+2ab+b^2,\qquad (a+b)^3=a^3+3a^2b+3ab^2+b^3.
\]

\begin{theorem}{הבינום של ניוטון}{binom}
לכל $a,b\in\R$ ולכל $n\ge1$:
\[
  (a+b)^n=\sum_{k=0}^{n}\binom{n}{k}a^k\,b^{n-k}
  = a^0b^n+\binom{n}{1}a^1b^{n-1}+\dots+\binom{n}{n-1}a^{n-1}b+\binom{n}{n}a^nb^0.
\]
\end{theorem}

\subsection{שימושים של הבינום של ניוטון}

\textbf{שימוש 1.} ניקח $a=b=1$:
\[
  (1+1)^n=2^n=\sum_{k=0}^{n}\binom{n}{k}\cdot1^k\cdot1^{n-k}=\sum_{k=0}^{n}\binom{n}{k}.
\]

\begin{example}{$n=10$}{}
\[
  2^{10}=\binom{10}{0}+\binom{10}{1}+\dots+\binom{10}{9}+\binom{10}{10}.
\]
כל מחובר סופר תתי-קבוצות של $\set{1,\dots,10}$ בגודל מסוים. למשל, $\binom{10}{0}=1$ סופר את תת-הקבוצה היחידה בגודל $0$, כלומר $\emptyset$. $\binom{10}{1}=10$ סופר את $\set{1},\set{2},\dots,\set{10}$, וכן הלאה.
\end{example}

\begin{keybox}
\[
  \binom{n}{k}=\text{\RL{מספר תתי-הקבוצות בגודל $k$ של קבוצה עם $n$ איברים.}}
\]
למשל, $\binom{4}{2}=6$ (ראו את טבלת הבחירה של 2 מתוך $\set{A,B,C,D}$).
\end{keybox}

\begin{corollary}{}{powsize}
אם ב-$A$ יש $n$ איברים, אז
\[
  \big|2^A\big|=\big|\Pow(A)\big|=2^n,
\]
כלומר ל-$A$ יש $2^n$ תתי-קבוצות. (נובע מהשימוש הראשון: סוכמים את מספר תתי-הקבוצות מכל גודל $k=0,1,\dots,n$.)
\end{corollary}

\begin{attention}{השימוש השני חסר}
כותרת הסעיף ברשימות היא "שני שימושים של הבינום של ניוטון", אבל מופיע רק שימוש אחד. ייתכן שהשני נכתב רק על הלוח או נלמד בהרצאה הבאה. יש להשלים.
\end{attention}

\begin{suggestion}{מועמד לשימוש השני}
שימוש נפוץ: ניקח $a=-1$, $b=1$. לכל $n\ge1$:
\[
  0=(-1+1)^n=\sum_{k=0}^{n}(-1)^k\binom{n}{k}
  \quad\Longrightarrow\quad
  \binom{n}{0}+\binom{n}{2}+\binom{n}{4}+\cdots=\binom{n}{1}+\binom{n}{3}+\cdots.
\]
במילים: לקבוצה עם $n\ge1$ איברים יש אותו מספר של תתי-קבוצות בגודל זוגי ובגודל אי-זוגי. למשל עבור $n=3$: $1+3=3+1$.
\end{suggestion}

\end{document}
